% =========================================================
%  config/packages.tex
%  Daftar paket LaTeX untuk dokumen RPS
%  Semua paket global didefinisikan di sini, bukan di sections/.
% =========================================================

% --- Pengkodean & bahasa --------------------------------
\usepackage[utf8]{inputenc}
\usepackage[T1]{fontenc}
\usepackage[english,indonesian]{babel}

% --- Tata letak halaman --------------------------------
\usepackage[a4paper,left=2.5cm,right=2.5cm,top=2.5cm,bottom=2.5cm]{geometry}
\usepackage{pdflscape}
\usepackage{fancyhdr}
\usepackage{ragged2e}

% --- Tabel ---------------------------------------------
% Catatan: opsi [table] pada xcolor (colortbl) tidak dipakai karena
% colortbl versi pada instalasi ini memerlukan array.sty yang lebih baru
% (\insert@pcolumn tidak terdefinisi). Pewarnaan baris tabel karena itu
% digantikan oleh penebalan garis header pada config/styles.tex.
\usepackage{booktabs}
\usepackage{array}
\usepackage{tabularx}
\usepackage{longtable}
\usepackage{multirow}
\usepackage{makecell}
\usepackage{xcolor}

% --- Gambar & posisi mengambang ------------------------
\usepackage{graphicx}
\usepackage{float}

% --- Tipografi & daftar --------------------------------
\usepackage{enumitem}
\usepackage{titlesec}
\usepackage{caption}

% --- Tautan (dimuat paling akhir) ----------------------
\usepackage[bookmarks=true]{hyperref}
\usepackage{url}
\usepackage{xurl}   % memungkinkan pemenggalan URL panjang pada sembarang karakter
